\documentclass[a4paper,11pt]{article}

\usepackage[margin=2.3cm]{geometry}
\usepackage{titlesec}
\usepackage{enumitem}
\usepackage[hidelinks]{hyperref}
\usepackage{xcolor}

\pagestyle{empty}
\setlength{\parindent}{0pt}
\setlength{\parskip}{6pt}

\titleformat{\section}
{\large\bfseries}
{\thesection}
{0.6em}
{}[\titlerule]

\titleformat{\subsection}
{\bfseries}
{\thesubsection}
{0.6em}
{}

\begin{document}

\begin{center}
	{\LARGE \textbf{Vision-Language-Action Models}}\\
	\vspace{2pt}
	Machine Learning and Deep Learning -- Presentation Notes\\
	\vspace{2pt}
	\textit{Assigned subtopics: RT-2, OpenVLA, Octo, $\pi_0$, Generalist Robot Policies, Vision-Language-Conditioned Control}
\end{center}

\vspace{6pt}

\section{The Big Picture: What a VLA Model Is}

A Vision-Language-Action (VLA) model is a single neural network that takes \textbf{(camera image, natural-language instruction)} as input and outputs a \textbf{robot action} directly -- with no separate, hand-engineered perception, planning, or control modules in between.

The core idea: take a model already pretrained on huge internet-scale image+text data (a Vision-Language Model -- the kind behind image captioning and visual question answering), and fine-tune it so that instead of outputting text, it outputs robot actions. The underlying bet is that the semantic and world knowledge baked into the VLM (what a ``banana'' looks like, what ``put X in Y'' means, spatial relationships, object affordances) transfers to robot control and generalizes far better than training a robot policy from scratch on a small set of demonstrations.

This addresses a longstanding robotics problem: classic robot policies only work for the exact tasks and objects they were trained on. VLA models aim to be \textbf{generalist} -- one set of weights that follows arbitrary language instructions across many tasks, objects, and sometimes even different robot embodiments.

\section{How the Six Subtopics Relate to Each Other}

These are not six independent topics -- they form one lineage, framed by two organizing concepts.

\begin{center}
\begin{minipage}{0.85\linewidth}
\centering
\textbf{``Generalist robot policies''} --- the \emph{goal / category} all of these belong to\\
\textbf{``Vision-language-conditioned control''} --- the \emph{mechanism} that makes it possible\\[8pt]
\begin{tabular}{cccc}
RT-2 & OpenVLA & Octo & $\pi_0$ \\
(closed, Google) & (open, 7B, OXE data) & (open, small, OXE data) & (open weights, flow matching)
\end{tabular}
\end{minipage}
\end{center}

\begin{itemize}
	\item \textbf{RT-2} established the paradigm (2023): fine-tune a VLM to speak ``action tokens'' instead of words.
	\item \textbf{OpenVLA} and \textbf{Octo} are both open-source responses trained on the same large public dataset (\textbf{Open X-Embodiment}, roughly 1M robot trajectories pooled across $\sim$20 labs and robot types), but they diverge architecturally in \emph{how} they turn language and vision into actions.
	\item \textbf{$\pi_0$} is a further evolution -- instead of discretizing actions into tokens (like RT-2/OpenVLA), it generates smooth, continuous, high-frequency actions via flow matching (related to diffusion models).
	\item \textbf{Generalist robot policies} is the umbrella term for what all four of these \emph{are} -- a shift away from one-task-one-policy toward one-model-many-tasks.
	\item \textbf{Vision-language-conditioned control} is \emph{how} they work internally -- the mechanism (attention over image patches and language tokens, jointly conditioning the action decoder) that each model implements differently.
\end{itemize}

A presentation covering all six bullet points naturally uses the last two as the framing/comparison narrative around the four concrete models.

\section{Each Subtopic in Detail}

\subsection{RT-2 (Robotic Transformer 2)}
Google DeepMind, 2023. Takes a pretrained VLM (PaLI-X or PaLM-E) and fine-tunes it \emph{jointly} on web-scale vision-language data and robot trajectory data. Robot actions (7 numbers: $x, y, z$, roll, pitch, yaw, gripper) are discretized into bins and represented as extra ``words'' in the model's vocabulary, so the model literally outputs a string of action-tokens the same way it would output a text answer. Because it is co-trained on web data, it can follow instructions involving objects/concepts never seen in robot demonstrations (e.g. ``pick up the extinct animal'' $\rightarrow$ picks the dinosaur toy), showing genuine semantic transfer. \textbf{Not open-sourced} -- no public weights.

\subsection{OpenVLA}
Stanford / Berkeley / Toyota Research / Google, 2024. An open-source response to RT-2. 7B parameters, built on a Llama-2 backbone fused with visual encoders (DINOv2 + SigLIP features). Trained on 970k trajectories from Open X-Embodiment. Same discrete-action-token approach as RT-2 (autoregressively predicts 256-bin-discretized action dimensions). Ships with reference inference code: feed one RGB image plus a text instruction, and it returns a 7-DoF action vector for one timestep. This is the standard reference point for ``what a VLA model actually is'' in current coursework.

\subsection{Octo}
UC Berkeley, 2024. Also trained on Open X-Embodiment, but architecturally different: a modular transformer where observations (multiple camera views, proprioception, goal images, or language) are tokenized flexibly and combined via cross-attention, and the action head is a \textbf{diffusion policy} that predicts a \emph{chunk} of future actions at once, rather than one action at a time autoregressively. Much smaller than OpenVLA -- 27M (``Octo-Small'') or 93M (``Octo-Base'') parameters, vs. 7B -- designed explicitly to be fast to fine-tune on a new robot with limited demonstrations. This is also the practical reason it is the easiest of the four to actually run live.

\subsection{$\pi_0$ (Pi-Zero)}
Physical Intelligence, 2024 (open weights released later via their \texttt{openpi} repository). Pushes further than discretized tokens (RT-2/OpenVLA) or per-chunk diffusion (Octo): it uses \textbf{flow matching} to generate continuous, high-frequency action trajectories directly, aimed at more dexterous tasks (e.g. folding laundry, plugging in a cable) that need smoother, faster control than token-by-token generation allows. Built on a VLM backbone (PaliGemma) with an added ``action expert'' module.

\subsection{Generalist Robot Policies (Concept)}
The category RT-2, OpenVLA, Octo, and $\pi_0$ all belong to: a single trained model that transfers across tasks, objects, scenes, and sometimes even different robot bodies (cross-embodiment), as opposed to the older paradigm of training one narrow policy per task from scratch.

\subsection{Vision-Language-Conditioned Control (Concept)}
The mechanism by which a model architecturally fuses a camera image and a language instruction into a shared representation that then conditions the action-generation process. This is the technical thread that differentiates the four models: discrete tokens with autoregressive decoding (RT-2/OpenVLA), a diffusion action head (Octo), and flow matching (\(\pi_0\)) -- three different answers to the same underlying question of how to turn (image, language) into a continuous robot action.

\section{What the Application / Demo Looks Like in Practice}

Regardless of which model is used, a working VLA system always runs the same control loop:

\begin{enumerate}
	\item A camera captures an RGB frame of the workspace.
	\item A natural-language command is provided, e.g. \textit{``pick up the red block and place it on the plate.''}
	\item Model forward pass: (image, instruction) $\rightarrow$ action.
	\item The action -- either a single 7-DoF vector, or a short chunk of future actions for Octo/$\pi_0$ -- is sent to the robot's controller, typically as a target end-effector delta-pose plus a gripper command.
	\item The loop repeats at some control frequency (a few Hz for token-based models like OpenVLA; higher for chunked/continuous models like Octo/$\pi_0$) until the task is judged complete.
\end{enumerate}

For OpenVLA and Octo specifically, this is exactly what their public example/inference code does out of the box: load the model, pass in an image and an instruction, and get an action vector back. Whether that action is then executed on a real robot arm, a simulated arm, or simply visualized/printed is an implementation decision independent of the model itself, and can be adapted to whatever hardware or simulation environment is available for the demo (real robot, simulation, or visualization-only).

\section{Demo Planning Note}

The presentation covers all six subtopics above. The live demo component may focus on a subset of the four concrete models (OpenVLA and Octo, or up to three including $\pi_0$), depending on what hardware/simulation access is confirmed and how reliably each model's public inference code can be set up in advance of the session.

\end{document}
